\documentclass{article}
\usepackage{amsmath,amssymb,amsthm}
\usepackage{algorithm}
\usepackage[noend]{algpseudocode}
\usepackage{fullpage}
\usepackage{hyperref}
\newtheorem{theorem}{Theorem}
\newtheorem{lemma}{Lemma}
\newtheorem{assumption}{Assumption}
\newcommand{\E}{\mathbb{E}}
\newcommand{\norm}[1]{\left\|#1\right\|}
\newcommand{\inner}[2]{\langle #1,#2\rangle}
\newcommand{\argmin}{\operatorname*{arg\,min}}

\begin{document}

\section*{Algorithm Name}
\textbf{SVRG (Stochastic Variance Reduced Gradient)} – a double–loop method for minimizing a finite sum
\[
f(x)=\frac1n\sum_{i=1}^{n}f_i(x),\qquad x\in\mathbb{R}^{d}.
\]

\section*{Mathematical Formulation}
Let $K\ge 1$ be the number of outer epochs and $m\ge 1$ the inner loop length.  
Given a stepsize $\eta>0$, the iterates are defined as follows.

\begin{enumerate}
\item \textbf{Outer loop (epoch $k$).}  Set the \emph{snapshot} point
\[
\tilde{x}^{k}=x_{0}^{k}\in\mathbb{R}^{d},
\qquad 
\tilde{\mu}^{k}\;:=\;\nabla f(\tilde{x}^{k})=\frac1n\sum_{i=1}^{n}\nabla f_i(\tilde{x}^{k}).
\]
\item \textbf{Inner loop ($t=0,\dots,m-1$).}  Sample an index $i_t$ uniformly from $\{1,\dots,n\}$ and compute the variance–reduced estimator
\[
v_t^{k}\;:=\;\nabla f_{i_t}\!\bigl(x_t^{k}\bigr)-\nabla f_{i_t}\!\bigl(\tilde{x}^{k}\bigr)+\tilde{\mu}^{k}.
\]
Update
\[
x_{t+1}^{k}=x_t^{k}-\eta\,v_t^{k}.
\]
\item After $m$ inner steps set the next snapshot
\[
\tilde{x}^{k+1}=x_{m}^{k}.
\]
\end{enumerate}
The algorithm returns $\tilde{x}^{K}$.

\section*{Pseudocode}
\begin{algorithm}[H]
\caption{SVRG}
\begin{algorithmic}[1]
\Require Number of epochs $K$, inner length $m$, stepsize $\eta>0$, initial point $x_0$.
\State $\tilde{x}^{0}\gets x_0$
\For{$k=0$ \textbf{to} $K-1$}
    \State $\tilde{\mu}^{k}\gets \frac1n\sum_{i=1}^{n}\nabla f_i(\tilde{x}^{k})$ \Comment{full gradient}
    \State $x_0^{k}\gets \tilde{x}^{k}$
    \For{$t=0$ \textbf{to} $m-1$}
        \State Sample $i_t\sim\text{Uniform}\{1,\dots,n\}$
        \State $v_t^{k}\gets \nabla f_{i_t}(x_t^{k})-\nabla f_{i_t}(\tilde{x}^{k})+\tilde{\mu}^{k}$
        \State $x_{t+1}^{k}\gets x_t^{k}-\eta\,v_t^{k}$
    \EndFor
    \State $\tilde{x}^{k+1}\gets x_{m}^{k}$
\EndFor
\State \Return $\tilde{x}^{K}$
\end{algorithmic}
\end{algorithm}

\section*{Dry Run Example}
Minimize $f(w)=(w-3)^{2}$, i.e. $n=1$, $f_1(w)= (w-3)^2$, $\nabla f_1(w)=2(w-3)$.  
Choose $x_0=0$, stepsize $\eta=0.1$, $m=3$, $K=1$.

\begin{center}
\begin{tabular}{c|c|c|c}
$t$ & $x_t$ & $v_t$ (since $i_t=1$) & $f(x_t)$\\ \hline
0 & $0$ & $2(0-3)-2(0-3)+2(0-3)=2(0-3)=-6$ & $9$\\
  & $x_1=0-0.1(-6)=0.6$ & $2(0.6-3)-2(0-3)+(-6)=2(-2.4)+6-6=-4.8$ & $5.76$\\
1 & $0.6$ & $v_1=-4.8$ & $5.76$\\
  & $x_2=0.6-0.1(-4.8)=1.08$ & $2(1.08-3)-2(0-3)+(-6)=2(-1.92)+6-6=-3.84$ & $3.6864$\\
2 & $1.08$ & $v_2=-3.84$ & $3.6864$\\
  & $x_3=1.08-0.1(-3.84)=1.464$ & – & $2.1410$\\
\end{tabular}
\end{center}
After the inner loop we set $\tilde{x}^{1}=x_3=1.464$, which is already closer to the optimum $w^{*}=3$.

\newpage
\section*{Assumptions}
\begin{assumption}[Finite sum]\label{assump:finite}
$f(x)=\frac1n\sum_{i=1}^{n}f_i(x)$ with each $f_i:\mathbb{R}^{d}\to\mathbb{R}$ differentiable.
\end{assumption}
\begin{assumption}[Smoothness]\label{assump:smooth}
Each $f_i$ is $L$–smooth, i.e.
\[
\norm{\nabla f_i(x)-\nabla f_i(y)}\le L\norm{x-y},\qquad\forall x,y\in\mathbb{R}^{d}.
\]
Consequently,
\[
f_i(y)\le f_i(x)+\inner{\nabla f_i(x)}{y-x}+\frac{L}{2}\norm{y-x}^{2}.
\]
\end{assumption}
\begin{assumption}[Strong convexity]\label{assump:strong}
$f$ is $\mu$–strongly convex ($0<\mu\le L$), i.e.
\[
f(y)\ge f(x)+\inner{\nabla f(x)}{y-x}+\frac{\mu}{2}\norm{y-x}^{2},\qquad\forall x,y.
\]
\end{assumption}
\begin{assumption}[Stepsize and inner length]\label{assump:params}
The stepsize satisfies $0<\eta<\frac{1}{2L}$ and the inner loop length $m$ is a positive integer (often $m=2n$ in practice).
\end{assumption}

\section*{Main Convergence Theorem}
\begin{theorem}[Linear convergence of SVRG]\label{thm:linear}
Under Assumptions~\ref{assump:finite}–\ref{assump:params}, let $\{ \tilde{x}^{k}\}_{k\ge0}$ be the sequence generated by SVRG. Define
\[
\rho\;:=\;\frac{1}{\mu\eta(1-2L\eta)m+1}\;\in(0,1).
\]
Then for every epoch $k\ge0$,
\[
\E\!\left[ f(\tilde{x}^{k+1})-f(x^{*})\right]\;\le\;\rho\;\E\!\left[ f(\tilde{x}^{k})-f(x^{*})\right],
\]
where $x^{*}=\argmin_{x}f(x)$ and the expectation is taken over all random draws of the indices $\{i_t\}$.
Consequently,
\[
\E\!\left[ f(\tilde{x}^{K})-f(x^{*})\right]\;\le\;\rho^{K}\,\bigl(f(\tilde{x}^{0})-f(x^{*})\bigr).
\]
\end{theorem}

\section*{Theoretical Properties}
\begin{itemize}
\item \textbf{Stability.} The contraction factor $\rho$ is decreasing in $\eta$ (up to $\frac1{2L}$) and in $m$; larger inner loops yield faster linear decay.
\item \textbf{Hyperparameter sensitivity.} If $\eta$ violates $\eta<\frac1{2L}$ the term $(1-2L\eta)$ becomes non‑positive and the bound collapses; thus the stepsize must respect the smoothness constant.
\item \textbf{Comparison.} Compared with vanilla SGD, which only achieves $O(1/T)$ sub‑linear rate for strongly convex problems, SVRG attains a geometric rate $O(\rho^{K})$ while still using stochastic gradients.
\item \textbf{Special case $m=n$.} When $m=n$, the bound matches the original SVRG analysis; setting $m=2n$ and $\eta=\frac{1}{10L}$ gives $\rho\approx\frac{2}{3}$.
\end{itemize}

\section*{Discussion}
The theorem shows that each outer epoch reduces the optimality gap by a constant factor $\rho<1$. The proof hinges on three elementary lemmas: smoothness, strong convexity, and a variance bound for the estimator $v_t^{k}$. The next sections present the full derivation with every algebraic step explicitly numbered.

\newpage
\section*{Preliminaries (Definitions and Lemmas)}
\begin{lemma}[Smoothness inequality]\label{lem:smooth}
For any $L$‑smooth function $h$,
\[
h(y)\le h(x)+\inner{\nabla h(x)}{y-x}+\frac{L}{2}\norm{y-x}^{2}.
\]
\end{lemma}
\begin{lemma}[Strong convexity inequality]\label{lem:strong}
For any $\mu$‑strongly convex $h$,
\[
\inner{\nabla h(x)}{x-x^{*}}\ge h(x)-h(x^{*})+\frac{\mu}{2}\norm{x-x^{*}}^{2}.
\]
\end{lemma}
\begin{lemma}[Variance bound of SVRG estimator]\label{lem:var}
For the estimator $v_t^{k}$ defined in the algorithm,
\[
\E\!\bigl[\|v_t^{k}-\nabla f(x_t^{k})\|^{2}\mid x_t^{k}\bigr]
\;\le\;2L\bigl(f(x_t^{k})-f(\tilde{x}^{k})- \inner{\nabla f(\tilde{x}^{k})}{x_t^{k}-\tilde{x}^{k}}\bigr).
\]
\end{lemma}
\begin{proof}[Proof of Lemma~\ref{lem:var}]
By definition,
\[
v_t^{k}-\nabla f(x_t^{k})
=
\bigl(\nabla f_{i_t}(x_t^{k})-\nabla f_{i_t}(\tilde{x}^{k})\bigr)
-\bigl(\nabla f(x_t^{k})-\nabla f(\tilde{x}^{k})\bigr).
\]
Take expectation over the uniform choice of $i_t$ and use independence:
\[
\E\!\bigl[\|v_t^{k}-\nabla f(x_t^{k})\|^{2}\mid x_t^{k}\bigr]
=
\E_{i_t}\!\Bigl[\bigl\|\nabla f_{i_t}(x_t^{k})-\nabla f_{i_t}(\tilde{x}^{k})
 -(\nabla f(x_t^{k})-\nabla f(\tilde{x}^{k}))\bigr\|^{2}\Bigr].
\]
Expanding the square and using $\E_{i_t}[\nabla f_{i_t}(\cdot)]=\nabla f(\cdot)$ gives
\[
\begin{aligned}
&\E_{i_t}\!\bigl[\|\nabla f_{i_t}(x_t^{k})-\nabla f_{i_t}(\tilde{x}^{k})\|^{2}\bigr]
-\|\nabla f(x_t^{k})-\nabla f(\tilde{x}^{k})\|^{2}\\
&\le\frac1n\sum_{i=1}^{n}\|\nabla f_i(x_t^{k})-\nabla f_i(\tilde{x}^{k})\|^{2}
\;\le\;2L\Bigl(f(x_t^{k})-f(\tilde{x}^{k})-\inner{\nabla f(\tilde{x}^{k})}{x_t^{k}-\tilde{x}^{k}}\Bigr),
\end{aligned}
\]
where the last inequality is a standard consequence of $L$‑smoothness (see e.g. Lemma~2 in the original SVRG paper). ∎
\end{proof}

\newpage
\section*{Main Proof (Step-by-Step Derivation)}
We prove Theorem~\ref{thm:linear}.  Throughout the proof, expectations are taken over all randomness up to the current iteration.

\bigskip
\noindent\textbf{Notation.} Let $x^{*}$ denote the unique minimizer of $f$.  For brevity write $x_t:=x_t^{k}$, $v_t:=v_t^{k}$, $\tilde{x}:=\tilde{x}^{k}$ and $\tilde{\mu}:=\nabla f(\tilde{x})$.

\bigskip
\noindent\textbf{Step 1. One–step distance recursion.}  
From the update $x_{t+1}=x_t-\eta v_t$ we have
\[
\begin{aligned}
\norm{x_{t+1}-x^{*}}^{2}
&=\norm{x_t-x^{*}}^{2}-2\eta\inner{v_t}{x_t-x^{*}}+\eta^{2}\norm{v_t}^{2}. \tag{1}
\end{aligned}
\]
Taking expectation conditioned on $x_t$ and using $\E[v_t\mid x_t]=\nabla f(x_t)$ (unbiasedness) yields
\[
\E\!\bigl[\norm{x_{t+1}-x^{*}}^{2}\mid x_t\bigr]
= \norm{x_t-x^{*}}^{2}
-2\eta\inner{\nabla f(x_t)}{x_t-x^{*}}
+\eta^{2}\E\!\bigl[\norm{v_t}^{2}\mid x_t\bigr]. \tag{2}
\]

\bigskip
\noindent\textbf{Step 2. Bound on $\E[\|v_t\|^{2}]$.}  
Decompose $\|v_t\|^{2}$ as
\[
\|v_t\|^{2}
= \|\nabla f(x_t)\|^{2}
+ \|v_t-\nabla f(x_t)\|^{2}
+2\inner{v_t-\nabla f(x_t)}{\nabla f(x_t)}.
\]
The cross term vanishes after conditioning because $\E[v_t-\nabla f(x_t)\mid x_t]=0$. Hence
\[
\E\!\bigl[\|v_t\|^{2}\mid x_t\bigr]
= \|\nabla f(x_t)\|^{2}
+ \E\!\bigl[\|v_t-\nabla f(x_t)\|^{2}\mid x_t\bigr]. \tag{3}
\]
Apply Lemma~\ref{lem:var} to bound the variance term:
\[
\E\!\bigl[\|v_t-\nabla f(x_t)\|^{2}\mid x_t\bigr]
\le 2L\Bigl(f(x_t)-f(\tilde{x})-\inner{\tilde{\mu}}{x_t-\tilde{x}}\Bigr). \tag{4}
\]

\bigskip
\noindent\textbf{Step 3. Use smoothness to bound $\|\nabla f(x_t)\|^{2}$.}  
From $L$‑smoothness and strong convexity we have
\[
\|\nabla f(x_t)\|^{2}
\le 2L\bigl(f(x_t)-f(x^{*})\bigr). \tag{5}
\]
(Proof: combine $L$‑smoothness inequality with $\nabla f(x^{*})=0$.)

\bigskip
\noindent\textbf{Step 4. Substitute (4)–(5) into (3).}  
\[
\E\!\bigl[\|v_t\|^{2}\mid x_t\bigr]
\le 2L\bigl(f(x_t)-f(x^{*})\bigr)
+2L\Bigl(f(x_t)-f(\tilde{x})-\inner{\tilde{\mu}}{x_t-\tilde{x}}\Bigr). \tag{6}
\]

\bigskip
\noindent\textbf{Step 5. Plug (6) into (2).}  
\[
\begin{aligned}
\E\!\bigl[\norm{x_{t+1}-x^{*}}^{2}\mid x_t\bigr]
\le &\;\norm{x_t-x^{*}}^{2}
-2\eta\inner{\nabla f(x_t)}{x_t-x^{*}}\\
&\;+\eta^{2}\,2L\bigl(f(x_t)-f(x^{*})\bigr)
+\eta^{2}\,2L\Bigl(f(x_t)-f(\tilde{x})-\inner{\tilde{\mu}}{x_t-\tilde{x}}\Bigr). \tag{7}
\end{aligned}
\]

\bigskip
\noindent\textbf{Step 6. Replace the inner product using strong convexity (Lemma~\ref{lem:strong}).}  
From Lemma~\ref{lem:strong} with $h=f$,
\[
\inner{\nabla f(x_t)}{x_t-x^{*}}
\ge f(x_t)-f(x^{*})+\frac{\mu}{2}\norm{x_t-x^{*}}^{2}. \tag{8}
\]
Insert (8) into (7):
\[
\begin{aligned}
\E\!\bigl[\norm{x_{t+1}-x^{*}}^{2}\mid x_t\bigr]
\le &\;\norm{x_t-x^{*}}^{2}
-2\eta\Bigl(f(x_t)-f(x^{*})+\frac{\mu}{2}\norm{x_t-x^{*}}^{2}\Bigr)\\
&\;+\eta^{2}2L\bigl(f(x_t)-f(x^{*})\bigr)
+\eta^{2}2L\Bigl(f(x_t)-f(\tilde{x})-\inner{\tilde{\mu}}{x_t-\tilde{x}}\Bigr). \tag{9}
\end{aligned}
\]

\bigskip
\noindent\textbf{Step 7. Collect terms w.r.t. $f(x_t)-f(x^{*})$.}  
Define $\Delta_t:=f(x_t)-f(x^{*})$ and $\Delta_{\tilde}:=f(\tilde{x})-f(x^{*})$.  Equation (9) becomes
\[
\begin{aligned}
\E\!\bigl[\norm{x_{t+1}-x^{*}}^{2}\mid x_t\bigr]
\le &\;\bigl(1-\mu\eta\bigr)\norm{x_t-x^{*}}^{2}\\
&\; -2\eta\Delta_t +2L\eta^{2}\Delta_t
+2L\eta^{2}\Bigl(\Delta_t-\Delta_{\tilde}-\inner{\tilde{\mu}}{x_t-\tilde{x}}\Bigr). \tag{10}
\end{aligned}
\]

\bigskip
\noindent\textbf{Step 8. Use the identity $\inner{\tilde{\mu}}{x_t-\tilde{x}} = f(\tilde{x})-f(x_t)+\underbrace{\bigl(\inner{\tilde{\mu}}{x_t-\tilde{x}}+f(x_t)-f(\tilde{x})\bigr)}_{\ge \frac{\mu}{2}\norm{x_t-\tilde{x}}^{2}}$ from strong convexity applied at $\tilde{x}$.  
Specifically,
\[
\inner{\tilde{\mu}}{x_t-\tilde{x}}
\ge f(x_t)-f(\tilde{x})+\frac{\mu}{2}\norm{x_t-\tilde{x}}^{2}. \tag{11}
\]
Thus
\[
-\inner{\tilde{\mu}}{x_t-\tilde{x}}
\le -\bigl(f(x_t)-f(\tilde{x})\bigr)-\frac{\mu}{2}\norm{x_t-\tilde{x}}^{2}. \tag{12}
\]

\bigskip
\noindent\textbf{Step 9. Substitute (12) into (10).}  
The term $2L\eta^{2}\bigl(-\inner{\tilde{\mu}}{x_t-\tilde{x}}\bigr)$ is bounded by
\[
2L\eta^{2}\Bigl[-\inner{\tilde{\mu}}{x_t-\tilde{x}}\Bigr]
\le 2L\eta^{2}\Bigl[-\bigl(f(x_t)-f(\tilde{x})\bigr)-\frac{\mu}{2}\norm{x_t-\tilde{x}}^{2}\Bigr]. \tag{13}
\]
Plugging (13) into (10) cancels the $+\bigl(f(x_t)-f(\tilde{x})\bigr)$ part, yielding
\[
\begin{aligned}
\E\!\bigl[\norm{x_{t+1}-x^{*}}^{2}\mid x_t\bigr]
\le &\;(1-\mu\eta)\norm{x_t-x^{*}}^{2}
-2\eta\Delta_t +2L\eta^{2}\Delta_t\\
&\; -2L\eta^{2}\frac{\mu}{2}\norm{x_t-\tilde{x}}^{2}. \tag{14}
\end{aligned}
\]

\bigskip
\noindent\textbf{Step 10. Drop the non‑positive term $-L\mu\eta^{2}\norm{x_t-\tilde{x}}^{2}$ to obtain a simpler recursion}
\[
\E\!\bigl[\norm{x_{t+1}-x^{*}}^{2}\mid x_t\bigr]
\le (1-\mu\eta)\norm{x_t-x^{*}}^{2}
-\bigl(2\eta-2L\eta^{2}\bigr)\Delta_t. \tag{15}
\]

\bigskip
\noindent\textbf{Step 11. Take total expectation and sum over the $m$ inner steps.}  
Let $\Phi_t:=\E\bigl[\norm{x_t-x^{*}}^{2}\bigr]$ and $\Delta_t:=\E\bigl[f(x_t)-f(x^{*})\bigr]$.  Taking expectation of (15) gives
\[
\Phi_{t+1}\le (1-\mu\eta)\Phi_t-(2\eta-2L\eta^{2})\Delta_t. \tag{16}
\]
Summing (16) for $t=0,\dots,m-1$ yields
\[
\Phi_{m}\le (1-\mu\eta)^{m}\Phi_{0}
-(2\eta-2L\eta^{2})\sum_{t=0}^{m-1}\Delta_t. \tag{17}
\]

\bigskip
\noindent\textbf{Step 12. Relate the average inner optimality gap to the snapshot gap.}  
Because $f$ is $\mu$‑strongly convex, $\Phi_{0}= \E\!\bigl[\norm{\tilde{x}-x^{*}}^{2}\bigr]\le \frac{2}{\mu}\bigl(f(\tilde{x})-f(x^{*})\bigr)=\frac{2}{\mu}\Delta_{\tilde}$.  
Moreover, by convexity of $f$,
\[
\frac1m\sum_{t=0}^{m-1}f(x_t)\;\ge\; f\!\bigl(\tilde{x}^{k+1}\bigr),
\]
so that
\[
\frac1m\sum_{t=0}^{m-1}\Delta_t\;\ge\;\E\!\bigl[f(\tilde{x}^{k+1})-f(x^{*})\bigr]=:\Delta_{\tilde}^{+}. \tag{18}
\]

\bigskip
\noindent\textbf{Step 13. Combine (17)–(18).}  
Using $\Phi_{m}\ge 0$ and (18) we obtain
\[
0\le (1-\mu\eta)^{m}\frac{2}{\mu}\Delta_{\tilde}
- (2\eta-2L\eta^{2})\,m\,\Delta_{\tilde}^{+}.
\]
Rearranging gives
\[
\Delta_{\tilde}^{+}\le
\frac{(1-\mu\eta)^{m}}{m\eta(1- L\eta)}\;\Delta_{\tilde}.
\]
Since $0<\eta<\frac1{2L}$, we have $1- L\eta\ge 1-2L\eta>0$.  Moreover,
\[
(1-\mu\eta)^{m}\le \frac{1}{1+\mu\eta m},
\]
which follows from the inequality $1-x\le \frac{1}{1+x}$ for $x\ge0$.  Consequently,
\[
\Delta_{\tilde}^{+}\le
\frac{1}{\mu\eta(1-2L\eta)m+1}\;\Delta_{\tilde}.
\]
Define
\[
\rho:=\frac{1}{\mu\eta(1-2L\eta)m+1}\in(0,1),
\]
which is exactly the contraction factor stated in Theorem~\ref{thm:linear}.  Hence
\[
\E\!\bigl[f(\tilde{x}^{k+1})-f(x^{*})\bigr]\le\rho\,
\E\!\bigl[f(\tilde{x}^{k})-f(x^{*})\bigr],
\]
establishing the linear convergence claim. ∎

\section*{Rate Derivation (With Constants)}
From Theorem~\ref{thm:linear},
\[
\E\!\bigl[f(\tilde{x}^{K})-f(x^{*})\bigr]
\le\rho^{K}\bigl(f(\tilde{x}^{0})-f(x^{*})\bigr),
\qquad
\rho=\frac{1}{1+\mu\eta(1-2L\eta)m}.
\]
If we choose the common SVRG parameters $\eta=\frac{1}{10L}$ and $m=2n$, then
\[
\rho
= \frac{1}{1+\frac{\mu}{10L}\bigl(1-\frac{1}{5}\bigr)2n}
= \frac{1}{1+\frac{9\mu}{50L}n}
\le \exp\!\Bigl(-\frac{9\mu}{50L}n\Bigr),
\]
showing an exponential decay with the total number of stochastic gradient evaluations $K\cdot m = 2nK$.

\section*{Special Cases}
\begin{itemize}
\item \textbf{Pure gradient descent.}  Setting $m=1$ and using the full gradient in every inner step collapses SVRG to deterministic gradient descent, and $\rho= \frac{1}{1+\mu\eta(1-2L\eta)}$, the classical linear rate.
\item \textbf{Large stepsize limit.}  If $\eta\to\frac{1}{2L}$, the factor $(1-2L\eta)$ tends to $0$, and $\rho\to \frac{1}{1+\mu\eta m}$, i.e. the dependence on $L$ disappears; however the condition $\eta<\frac{1}{2L}$ must still hold.
\item \textbf{Infinite inner loop $m\to\infty$.}  The contraction factor tends to $0$, meaning the snapshot would converge in a single outer epoch (the algorithm would then be equivalent to full gradient descent).
\end{itemize}

\end{document}